\nonstopmode{}
\documentclass[letterpaper]{book}
\usepackage[times,inconsolata,hyper]{Rd}
\usepackage{makeidx}
\usepackage[utf8]{inputenc} % @SET ENCODING@
% \usepackage{graphicx} % @USE GRAPHICX@
\makeindex{}
\begin{document}
\chapter*{}
\begin{center}
{\textbf{\huge Package `rcicr'}}
\par\bigskip{\large \today}
\end{center}
\inputencoding{utf8}
\ifthenelse{\boolean{Rd@use@hyper}}{\hypersetup{pdftitle = {rcicr: Reverse-Correlation Image-Classification Toolbox}}}{}
\ifthenelse{\boolean{Rd@use@hyper}}{\hypersetup{pdfauthor = {Ron Dotsch}}}{}
\begin{description}
\raggedright{}
\item[Type]\AsIs{Package}
\item[Title]\AsIs{Reverse-Correlation Image-Classification Toolbox}
\item[Version]\AsIs{1.0.1.9000}
\item[URL]\AsIs{}\url{https://medium.com/@rondotsch/reverse-correlation-image-classification-using-r-a0701648fb0/}\AsIs{}
\item[Description]\AsIs{Functions to generate stimuli and analyze data of reverse
correlation image classification experiments (psychophysical tasks aimed at
visualizing cognitive mental representations of faces).}
\item[License]\AsIs{GPL-2}
\item[Encoding]\AsIs{UTF-8}
\item[Depends]\AsIs{R (>= 4.1)}
\item[Imports]\AsIs{matlab, png, jpeg, dplyr, scales, viridis, utils, parallel,
doParallel, foreach, spatstat.explore, spatstat.geom, raster,
tibble, yesno}
\item[Suggests]\AsIs{testthat (>= 3.0.0), covr, withr, knitr, rmarkdown}
\item[Config/testthat/edition]\AsIs{3}
\item[VignetteBuilder]\AsIs{knitr}
\item[RoxygenNote]\AsIs{7.3.1}
\item[NeedsCompilation]\AsIs{no}
\item[Author]\AsIs{Ron Dotsch [aut, cre]}
\item[Maintainer]\AsIs{Ron Dotsch }\email{rdotsch@gmail.com}\AsIs{}
\end{description}
\Rdcontents{\R{} topics documented:}
\inputencoding{utf8}
\HeaderA{rcicr-package}{Reverse-Correlation Image-Classification Toolbox}{rcicr.Rdash.package}
\aliasA{rcicr}{rcicr-package}{rcicr}
\keyword{package}{rcicr-package}
%
\begin{Description}
Toolbox with functions to generate stimuli and analyze data of reverse correlation image classification experiments. Reverse correlation is a psychophysical technique originally derived from signal detection theory. This package focuses on visualizing internal representations of participants using visual stimuli in a perceptual taks.
\end{Description}
%
\begin{Details}

\Tabular{ll}{
Package: & rcicr\\{}
Type: & Package\\{}
Version: & 0.4.0\\{}
Date: & 2017-07-25\\{}
License: & GPL-2\\{}
}
\bold{Generating stimuli}

Load the package with \code{library(rcicr)}. Then generate stimuli with:

\code{generateStimuli2IFC(base\_face\_files, n\_trials = 770)}

This will generate stimuli for 770 trials of a 2 images forced choice reverse correlation image classification task with sinusoid noise. By default the stimuli will have a resolution of 512 x 512 pixels. The stimuli will be saved as jpegs to a folder called stimuli in your current working directory, and an .Rdata file will be saved that contains the stimulus parameters necessary for analysis.

The \code{base\_face\_files} argument is a list of jpegs that should be used as base images for the stimuli. The base\_face\_files variable might look like this:

\code{base\_face\_files <- list('male'='male.jpg', 'female'='female.jpg')}

For each jpeg a set of stimuli will be created using the same noise patterns as for the other sets. The jpeg should have the resolution that you want the stimuli to have. By default this should be 512 x 512 pixels. If you want a different size, resize your base image to either 128 x 128 or 256 x 256 for smaller stimuli, or 1024 x 1024 for bigger stimuli. In that case, also set the \code{img\_size} parameter accordingly.

You are now ready to collect data with the stimuli you just created. The stimuli are named according to their sequence number when generating and whether the original noise is superimposed or the negative/inverted noise. Stimuli with the same sequence number should be presented side by side in the same trial. Record which stimulus a participant selected at any given trial (the original, or the inverted). At the very least be sure that in your data file the connection can be made between the response key of the participant and which stimulus was selected on each trial. Use any presentation software you like (I recommend python-based open source alternatives, like PsychoPy, Expyriment, or OpenSesame).

\bold{Data analysis}

Analyzing reverse correlation data is all about computing classification images. Use the following function for your data collected using the 2 images forced choice stimuli:

\code{ci <- generateCI2IFC(stimuli, responses, baseimage, rdata)}

The \code{stimuli} paramater should be a vector containing the sequence numbers of the stimuli that were presented in the trials of the task. The \code{responses} parameter contains, in the order of the stimuli vector, the response of a participant to those stimuli (coded 1 if the original stimulus was selected and -1 if the inverted stimulus was selected). The \code{baseimage} parameter is a string specifying which base image was used (not the file name, but the name in the list of \code{base\_face\_files}. So for the stimuli generated above, either \code{'male'} or \code{'female'}, depending on which set of stimuli was presented to the participant whose data you're analyzing). Finally, rdata is a string pointing to the .RData file that was created automatically when you generated the stimuli. It contains the parameters for each stimulus, necessary to create the classification image.

By default JPEGs of the classification images will be saved automatically. The returned values can be used later to optimally rescale the noise relative to the base image. For instance, if you have a list of cis from various participants (i.e., a list of the values returned by several calls to \code{generateCI2IFC}, one for each participant), you can use the autoscale function to generate classification images that are scaled identically and therefore straightforward to compare:

\code{scaled\_cis <- autoscale(cis, saveasjpegs = TRUE)}

\bold{Computing CIs for many participants or conditions}

Data analysis as described above can be automatized for a batch of participants or conditions using \code{batchGenerateCI2IFC}. Please see instructions for that function.

\bold{Note}

Currently, the package is still in alpha stage. Much may still change. It only supports 2 Image Forced Choice tasks, although the underlying functions can be used for other versions of the reverse correlation task. It also only supports sinusoid noise. In the future, it will support Gaussian white noise, as well as additional variants of the task.

If you use this package for your experiments, please cite the package in your publications. Use \code{citation('rcicr')} to print the appropriate citation for the current version of the package.
\end{Details}
%
\begin{Author}
Ron Dotsch <rdotsch@gmail.com> (http://ron.dotsch.org/)
Maintainer: Ron Dotsch <rdotsch@gmail.com>
\end{Author}
%
\begin{References}
Dotsch, R., \& Todorov, A. (2012). Reverse correlating social face perception. Social Psychological and Personality Science, 3 (5), 562-571.

Dotsch, R., Wigboldus, D. H. J., Langner, O., \& Van Knippenberg, A. (2008). Ethnic out-group faces are biased in the prejudiced mind. Psychological Science, 19, 978-980.
\end{References}
%
\begin{Examples}
\begin{ExampleCode}
#simple examples will be added soon.
\end{ExampleCode}
\end{Examples}
\inputencoding{utf8}
\HeaderA{autoscale}{Determines optimal scaling constant for a list of CIs}{autoscale}
%
\begin{Description}
Determines optimal scaling constant for a list of CIs
\end{Description}
%
\begin{Usage}
\begin{verbatim}
autoscale(cis, save_as_pngs = TRUE, targetpath = "./cis")
\end{verbatim}
\end{Usage}
%
\begin{Arguments}
\begin{ldescription}
\item[\code{cis}] List of cis, each of which are a list containing the pixel matrices of at least the noise pattern (\code{\$ci}) and if the noise patterns need to be written to PNGs, also the base image (\code{\$base}).

\item[\code{save\_as\_pngs}] Boolean, when set to true, the autoscaled noise patterns will be combined with their respective base images and saved as PNGs (using the key of the list as name).

\item[\code{targetpath}] Optional string specifying path to save PNGs to (default: ./cis).
\end{ldescription}
\end{Arguments}
%
\begin{Value}
The input \code{cis} list, with a \code{\$scaled} pixel matrix added to each element. The
scaling constant that was determined is printed to the console, not returned.
\end{Value}
%
\begin{Examples}
\begin{ExampleCode}
cis <- list(
  participant1 = list(ci = matrix(runif(64, -0.2, 0.2), 8, 8), base = matrix(0.5, 8, 8)),
  participant2 = list(ci = matrix(runif(64, -0.3, 0.3), 8, 8), base = matrix(0.5, 8, 8))
)
scaled_cis <- autoscale(cis, save_as_pngs = FALSE)
\end{ExampleCode}
\end{Examples}
\inputencoding{utf8}
\HeaderA{batchGenerateCI}{Generates multiple classification images by participant or condition}{batchGenerateCI}
%
\begin{Description}
Generate classification image for any reverse correlation task that displays independently generated alternatives.
\end{Description}
%
\begin{Usage}
\begin{verbatim}
batchGenerateCI(
  data,
  by,
  stimuli,
  responses,
  baseimage,
  rdata,
  save_as_png = TRUE,
  targetpath = "./cis",
  label = "",
  antiCI = FALSE,
  scaling = "autoscale",
  constant = 0.1
)
\end{verbatim}
\end{Usage}
%
\begin{Arguments}
\begin{ldescription}
\item[\code{data}] Data frame

\item[\code{by}] String specifying column name that specifies the smallest unit (participant, condition) to subset the data on and calculate CIs for.

\item[\code{stimuli}] String specifying column name in data frame that contains the stimulus numbers of the presented stimuli.

\item[\code{responses}] String specifying column name in data frame that contains the responses coded 1 for original stimulus selected and -1 for inverted stimulus selected.

\item[\code{baseimage}] String specifying which base image was used. Not the file name, but the key used in the list of base images at time of generating the stimuli.

\item[\code{rdata}] String pointing to .RData file that was created when stimuli were generated. This file contains the contrast parameters of all generated stimuli.

\item[\code{save\_as\_png}] Boolean stating whether to additionally save the CI as PNG image.

\item[\code{targetpath}] Optional string specifying path to save PNGs to (default: ./cis).

\item[\code{label}] Optional string to insert in file names of PNGs to make them easier to identify.

\item[\code{antiCI}] Optional boolean specifying whether antiCI instead of CI should be computed.

\item[\code{scaling}] Optional string specifying scaling method: \code{none}, \code{constant},  \code{independent} or \code{autoscale} (default).

\item[\code{constant}] Optional number specifying the value used as constant scaling factor for the noise (only works for \code{scaling='constant'}).
\end{ldescription}
\end{Arguments}
%
\begin{Details}
This function saves the classification images by participant or condition as PNG to a folder and returns the CIs.
\end{Details}
%
\begin{Value}
List of classification image data structures (which are themselves lists of pixel matrix of classification noise only, scaled classification noise only, base image only and combined).
\end{Value}
%
\begin{Examples}
\begin{ExampleCode}

# a synthetic square grayscale image stands in for a real base face photo
base_face <- tempfile(fileext = ".png")
png::writePNG(matrix(runif(32 * 32), 32, 32), base_face)

stimulus_path <- tempdir()
generateStimuli2IFC(
  base_face_files = list(face = base_face),
  n_trials = 6,
  img_size = 32,
  stimulus_path = stimulus_path,
  seed = 1,
  ncores = 1,
  nscales = 1,
  save_as_png = FALSE
)
rdata_file <- list.files(stimulus_path, pattern = "\\.Rdata$", full.names = TRUE)[1]

# two "participants", three trials each
data <- data.frame(
  participant = rep(c("p1", "p2"), each = 3),
  stimulus = 1:6,
  response = sample(c(1, -1), 6, replace = TRUE)
)

cis <- suppressWarnings(batchGenerateCI(
  data = data, by = "participant", stimuli = "stimulus", responses = "response",
  baseimage = "face", rdata = rdata_file, save_as_png = FALSE
))

\end{ExampleCode}
\end{Examples}
\inputencoding{utf8}
\HeaderA{batchGenerateCI2IFC}{Generates multiple 2IFC classification images by participant or condition}{batchGenerateCI2IFC}
%
\begin{Description}
Generate classification image for 2 images forced choice reverse correlation task.
\end{Description}
%
\begin{Usage}
\begin{verbatim}
batchGenerateCI2IFC(
  data,
  by,
  stimuli,
  responses,
  baseimage,
  rdata,
  save_as_png = TRUE,
  targetpath = "./cis",
  antiCI = FALSE,
  scaling = "autoscale",
  constant = 0.1,
  label = ""
)
\end{verbatim}
\end{Usage}
%
\begin{Arguments}
\begin{ldescription}
\item[\code{data}] Data frame

\item[\code{by}] String specifying column name that specifies the smallest unit (participant, condition) to subset the data on and calculate CIs for.

\item[\code{stimuli}] String specifying column name in data frame that contains the stimulus numbers of the presented stimuli.

\item[\code{responses}] String specifying column name in data frame that contains the responses coded 1 for original stimulus selected and -1 for inverted stimulus selected.

\item[\code{baseimage}] String specifying which base image was used. Not the file name, but the key used in the list of base images at time of generating the stimuli.

\item[\code{rdata}] String pointing to .RData file that was created when stimuli were generated. This file contains the contrast parameters of all generated stimuli.

\item[\code{save\_as\_png}] Boolean stating whether to additionally save the CI as PNG image.

\item[\code{targetpath}] Optional string specifying path to save PNGs to (default: ./cis).

\item[\code{antiCI}] Optional boolean specifying whether antiCI instead of CI should be computed.

\item[\code{scaling}] Optional string specifying scaling method: \code{none}, \code{constant}, \code{matched}, \code{independent}, or \code{autoscale} (default).

\item[\code{constant}] Optional number specifying the value used as constant scaling factor for the noise (only works for \code{scaling='constant'}).

\item[\code{label}] Optional string to insert in file names of PNGs to make them easier to identify.
\end{ldescription}
\end{Arguments}
%
\begin{Details}
This function saves the classification images by participant or condition as PNG to a folder and returns the CIs.
\end{Details}
%
\begin{Value}
List of classification image data structures (which are themselves lists of pixel matrix of classification noise only, scaled classification noise only, base image only and combined).
\end{Value}
%
\begin{Examples}
\begin{ExampleCode}

# a synthetic square grayscale image stands in for a real base face photo
base_face <- tempfile(fileext = ".png")
png::writePNG(matrix(runif(32 * 32), 32, 32), base_face)

stimulus_path <- tempdir()
generateStimuli2IFC(
  base_face_files = list(face = base_face),
  n_trials = 6,
  img_size = 32,
  stimulus_path = stimulus_path,
  seed = 1,
  ncores = 1,
  nscales = 1,
  save_as_png = FALSE
)
rdata_file <- list.files(stimulus_path, pattern = "\\.Rdata$", full.names = TRUE)[1]

# two "participants", three trials each
data <- data.frame(
  participant = rep(c("p1", "p2"), each = 3),
  stimulus = 1:6,
  response = sample(c(1, -1), 6, replace = TRUE)
)

cis <- suppressWarnings(batchGenerateCI2IFC(
  data = data, by = "participant", stimuli = "stimulus", responses = "response",
  baseimage = "face", rdata = rdata_file, save_as_png = FALSE
))

\end{ExampleCode}
\end{Examples}
\inputencoding{utf8}
\HeaderA{computeCumulativeCICorrelation}{Computes cumulative trial CIs correlations with final/target CI}{computeCumulativeCICorrelation}
%
\begin{Description}
Computes cumulative trial CIs correlations with final/target CI.
\end{Description}
%
\begin{Usage}
\begin{verbatim}
computeCumulativeCICorrelation(
  stimuli,
  responses,
  baseimage,
  rdata,
  targetci = list(),
  step = 1
)
\end{verbatim}
\end{Usage}
%
\begin{Arguments}
\begin{ldescription}
\item[\code{stimuli}] Vector with stimulus numbers (should be numeric) that were presented in the order of the response vector. Stimulus numbers must match those in file name of the generated stimuli.

\item[\code{responses}] Vector specifying the responses in the same order of the stimuli vector, coded 1 for original stimulus selected and -1 for inverted stimulus selected.

\item[\code{baseimage}] String specifying which base image was used. Not the file name, but the key used in the list of base images at time of generating the stimuli.

\item[\code{rdata}] String pointing to .RData file that was created when stimuli were generated. This file contains the contrast parameters of all generated stimuli.

\item[\code{targetci}] List Target CI object generated with rcicr functions to correlate cumulative CIs with.

\item[\code{step}] Step size in sequence of trials to compute correlations with.
\end{ldescription}
\end{Arguments}
%
\begin{Details}
Use for instance for plotting curves of trial-final/target CI correlations to estimate how many trials are necessary in your task
\end{Details}
%
\begin{Value}
Vector containing correlation between cumulative CI and final/target CI.
\end{Value}
%
\begin{Examples}
\begin{ExampleCode}

# a synthetic square grayscale image stands in for a real base face photo
base_face <- tempfile(fileext = ".png")
png::writePNG(matrix(runif(32 * 32), 32, 32), base_face)

stimulus_path <- tempdir()
generateStimuli2IFC(
  base_face_files = list(face = base_face),
  n_trials = 6,
  img_size = 32,
  stimulus_path = stimulus_path,
  seed = 1,
  ncores = 1,
  nscales = 1,
  save_as_png = FALSE
)
rdata_file <- list.files(stimulus_path, pattern = "\\.Rdata$", full.names = TRUE)[1]

responses <- sample(c(1, -1), 6, replace = TRUE)
correlations <- suppressWarnings(computeCumulativeCICorrelation(
  stimuli = 1:6, responses = responses, baseimage = "face", rdata = rdata_file
))

\end{ExampleCode}
\end{Examples}
\inputencoding{utf8}
\HeaderA{computeInfoVal2IFC}{Computes Informational Value}{computeInfoVal2IFC}
%
\begin{Description}
Computes Informational Value for a single CI in a 2IFC task.
\end{Description}
%
\begin{Usage}
\begin{verbatim}
computeInfoVal2IFC(target_ci, rdata, iter = 10000, force_gen_ref_dist = FALSE)
\end{verbatim}
\end{Usage}
%
\begin{Arguments}
\begin{ldescription}
\item[\code{target\_ci}] A classification image object (list-type) as returned by generateCI

\item[\code{rdata}] String pointing to .RData file that was created when stimuli were generated. This file contains the contrast parameters of all generated stimuli and possibly its corresponding reference distribution generated with generateReferenceDistribution().

\item[\code{iter}] Number of iterations for the simulation of the reference distribution (only used if reference distribution is not already pre-generated and present in rdata file)

\item[\code{force\_gen\_ref\_dist}] Boolean specifying whether to override the default behavior to use pre-computed values for the reference distribution for specific task parameters and instead force to recompute the reference distribution (default: FALSE).
\end{ldescription}
\end{Arguments}
%
\begin{Details}
The Informational Value metric can be considered as a z-score that quantifies the signal
present in a classification image. The higher the Informational Value, the more signal. It is
possible to use a cut-off such as z = 1.96 to select classification images with significant
signal under alpha = 0.05.

Informational Value is computed by simulating random responding under identical task parameters to
an empirical dataset (called the reference distribution). The metric quantifies how unlikely it is
to observe these data under the null-hypothesis that there is no signal (i.e., that there is only random responding).

The simulation to compute the reference distribution takes a long time, and is only run locally when
pre-computed values for the reference distribution matching the stimulus set in the .Rdata file have
not been supplied by the rcicr package.

For more information see Brinkman, Goffin, Aarts, van Haren, \& Dotsch (in prep).
\end{Details}
%
\begin{Value}
Informational value (z-score)
\end{Value}
%
\begin{Examples}
\begin{ExampleCode}

# a synthetic square grayscale image stands in for a real base face photo
base_face <- tempfile(fileext = ".png")
png::writePNG(matrix(runif(32 * 32), 32, 32), base_face)

stimulus_path <- tempdir()
generateStimuli2IFC(
  base_face_files = list(face = base_face),
  n_trials = 6,
  img_size = 32,
  stimulus_path = stimulus_path,
  seed = 1,
  ncores = 1,
  nscales = 1,
  save_as_png = FALSE
)
rdata_file <- list.files(stimulus_path, pattern = "\\.Rdata$", full.names = TRUE)[1]

# compute (and cache in rdata_file) a reference distribution; iter is kept
# tiny here for a fast example, in practice use iter >= 10000.
# Run from a temp working directory: generateReferenceDistribution2IFC()
# re-generates stimuli via generateStimuli2IFC() without forwarding a
# stimulus_path, so it always creates a ./stimuli directory relative to the
# working directory.
withr::with_dir(tempdir(), {
  suppressWarnings(generateReferenceDistribution2IFC(rdata_file, iter = 3, ncores = 1))
})

responses <- sample(c(1, -1), 6, replace = TRUE)
target_ci <- generateCI(
  stimuli = 1:6, responses = responses, baseimage = "face",
  rdata = rdata_file, save_as_png = FALSE
)

computeInfoVal2IFC(target_ci = target_ci, rdata = rdata_file)

\end{ExampleCode}
\end{Examples}
\inputencoding{utf8}
\HeaderA{deg2rad}{Convert angle in degrees to radians}{deg2rad}
%
\begin{Description}
Convert angle in degrees to radians
\end{Description}
%
\begin{Usage}
\begin{verbatim}
deg2rad(deg)
\end{verbatim}
\end{Usage}
%
\begin{Arguments}
\begin{ldescription}
\item[\code{deg}] Angle in degrees
\end{ldescription}
\end{Arguments}
%
\begin{Value}
The angle in radians.
\end{Value}
%
\begin{Examples}
\begin{ExampleCode}
deg2rad(180)
\end{ExampleCode}
\end{Examples}
\inputencoding{utf8}
\HeaderA{generateCI}{Generates classification image}{generateCI}
%
\begin{Description}
Generate classification image for any reverse correlation task.
\end{Description}
%
\begin{Usage}
\begin{verbatim}
generateCI(
  stimuli,
  responses,
  baseimage,
  rdata,
  participants = NA,
  save_individual_cis = FALSE,
  save_as_png = TRUE,
  filename = "",
  targetpath = "./cis",
  antiCI = FALSE,
  scaling = "independent",
  scaling_constant = 0.1,
  individual_scaling = "independent",
  individual_scaling_constant = 0.1,
  zmap = F,
  zmapmethod = "quick",
  zmapdecoration = T,
  sigma = 3,
  threshold = 3,
  zmaptargetpath = "./zmaps",
  n_cores = detectCores() - 1,
  mask = NA
)
\end{verbatim}
\end{Usage}
%
\begin{Arguments}
\begin{ldescription}
\item[\code{stimuli}] Vector with stimulus numbers (should be numeric) that were presented in the order of the response vector. Stimulus numbers must match those in file name of the generated stimuli.

\item[\code{responses}] Vector specifying the responses in the same order of the stimuli vector, coded 1 for original stimulus selected and -1 for inverted stimulus selected.

\item[\code{baseimage}] String specifying which base image was used. Not the file name, but the key used in the list of base images at time of generating the stimuli.

\item[\code{rdata}] String pointing to .RData file that was created when stimuli were generated. This file contains the contrast parameters of all generated stimuli.

\item[\code{participants}] Optional vector specifying participant IDs. If specified, will compute the requested CIs in two steps: step 1, compute CI for each participant. Step 2, compute final CI by averaging participant CIs. If unspecified, the function defaults to averaging all data in the stimuli and responses vector.

\item[\code{save\_individual\_cis}] Optional boolean specifying whether individual CIs should be save as PNG images when the \code{participants} parameter is used.

\item[\code{save\_as\_png}] Optional boolean stating whether to additionally save the CI as PNG image.

\item[\code{filename}] Optional string to specify a file name for the PNG image.

\item[\code{targetpath}] Optional string specifying path to save PNGs to (default: ./cis).

\item[\code{antiCI}] Optional boolean specifying whether antiCI instead of CI should be computed.

\item[\code{scaling}] Optional string specifying scaling method: \code{none}, \code{constant}, \code{matched}, or \code{independent} (default). This scaling applies to the group-level CIs if both individual-level and group-level CIs are being generated.

\item[\code{scaling\_constant}] Optional number specifying the value used as constant scaling factor for the noise (only works for \code{scaling='constant'}). This scaling applies to the group-level CIs if both individual-level and group-level CIs are being generated.

\item[\code{individual\_scaling}] Optional string specifying scaling method for individual CIs: \code{none}, \code{constant}, \code{independent} (default).

\item[\code{individual\_scaling\_constant}] Optional number specifying the value used as constant scaling factor for the noise of individual CIs (only works for \code{individual\_scaling='constant'}).

\item[\code{zmap}] Boolean specifying whether a z-map should be created (default: FALSE).

\item[\code{zmapmethod}] String specifying the method to create the z-map. Can be: \code{quick} (default), \code{t.test}.

\item[\code{zmapdecoration}] Optional boolean specifying whether the Z-map should be plotted with margins, text (sigma, threshold) and a scale (default: TRUE).

\item[\code{sigma}] Integer specifying the amount of smoothing to apply when generating the z-maps (default: 3).

\item[\code{threshold}] Integer specifying the threshold z-score (default: 3). Z-scores below the threshold will not be plotted on the z-map.

\item[\code{zmaptargetpath}] Optional string specifying path to save z-map PNGs to (default: ./zmaps).

\item[\code{n\_cores}] Optional integer specifying the number of CPU cores to use to generate the z-map (default: detectCores()-1).

\item[\code{mask}] Optional 2D matrix that defines the mask to be applied to the CI (1 = masked, 0 = unmasked). May also be a string specifying the path to a grayscale PNG image (black = masked, white = unmasked). Default: NA.
\end{ldescription}
\end{Arguments}
%
\begin{Details}
This function saves the classification image as PNG to a folder and returns the CI. Your choice of scaling
matters. The default, \code{'independent'}, picks the lowest scaling constant that avoids clipping this
particular classification image (see \code{'constant'} scaling below for the formula), so it is not
comparable across classification images with different noise ranges.

\code{'matched'} scaling will match the range of the intensity of the pixels to
the range of the base image pixels. This scaling is nonlinear and depends on the range of both base image
and noise pattern. It is truly suboptimal, because it shifts the 0 point of the noise (that is, pixels that would
not have changed the base image at all before scaling may change the base image after scaling and vice versa). It is
however the quick and dirty way to see how the CI noise affects the base image.

For more control, use \code{'constant'} scaling, where the scaling is independent of
the base image and noise range, but where the choice of constant is arbitrary (provided by the user with
the \code{constant} parameter). The noise is then scale as follows: \code{scaled <- (ci + constant) / (2*constant)}.
Note that pixels can take intensity values between 0 and 1. If your scaled noise exceeds those values,
a warning will be given. You should pick a higher constant (but do so consistently for different classification images
that you want to compare). The higher the constant, the less visible the noise will be in the resulting image.

When creating multiple classification images a good strategy is to find the lowest constant that works for all
classification images. This can be automatized using the \code{autoscale} function.
\end{Details}
%
\begin{Value}
List of pixel matrix of classification noise only, scaled classification noise only, base image only and combined.
\end{Value}
%
\begin{Examples}
\begin{ExampleCode}

# a synthetic square grayscale image stands in for a real base face photo
base_face <- tempfile(fileext = ".png")
png::writePNG(matrix(runif(32 * 32), 32, 32), base_face)

stimulus_path <- tempdir()
generateStimuli2IFC(
  base_face_files = list(face = base_face),
  n_trials = 6,
  img_size = 32,
  stimulus_path = stimulus_path,
  seed = 1,
  ncores = 1,
  nscales = 1,
  save_as_png = FALSE
)
rdata_file <- list.files(stimulus_path, pattern = "\\.Rdata$", full.names = TRUE)[1]

responses <- sample(c(1, -1), 6, replace = TRUE)
ci <- generateCI(
  stimuli = 1:6, responses = responses, baseimage = "face",
  rdata = rdata_file, save_as_png = FALSE
)

\end{ExampleCode}
\end{Examples}
\inputencoding{utf8}
\HeaderA{generateCI2IFC}{Generates 2IFC classification image}{generateCI2IFC}
%
\begin{Description}
Generate classification image for 2 images forced choice reverse correlation task.  This function exists for backwards compatibility. You can also just use \code{generateCI()}, which this function wraps.
\end{Description}
%
\begin{Usage}
\begin{verbatim}
generateCI2IFC(
  stimuli,
  responses,
  baseimage,
  rdata,
  save_as_png = TRUE,
  filename = "",
  targetpath = "./cis",
  antiCI = FALSE,
  scaling = "independent",
  constant = 0.1
)
\end{verbatim}
\end{Usage}
%
\begin{Arguments}
\begin{ldescription}
\item[\code{stimuli}] Vector with stimulus numbers (should be numeric) that were presented in the order of the response vector. Stimulus numbers must match those in file name of the generated stimuli.

\item[\code{responses}] Vector specifying the responses in the same order of the stimuli vector, coded 1 for original stimulus selected and -1 for inverted stimulus selected.

\item[\code{baseimage}] String specifying which base image was used. Not the file name, but the key used in the list of base images at time of generating the stimuli.

\item[\code{rdata}] String pointing to .RData file that was created when stimuli were generated. This file contains the contrast parameters of all generated stimuli.

\item[\code{save\_as\_png}] Boolean stating whether to additionally save the CI as PNG image.

\item[\code{filename}] Optional string to specify a file name for the PNG image.

\item[\code{targetpath}] Optional string specifying path to save PNGs to (default: ./cis).

\item[\code{antiCI}] Optional boolean specifying whether antiCI instead of CI should be computed.

\item[\code{scaling}] Optional string specifying scaling method: \code{none}, \code{constant}, \code{matched}, or \code{independent} (default).

\item[\code{constant}] Optional number specifying the value used as constant scaling factor for the noise (only works for \code{scaling='constant'}).
\end{ldescription}
\end{Arguments}
%
\begin{Details}
This function saves the classification image as PNG to a folder and returns the CI. Your choice of scaling
matters. The default, \code{'independent'}, picks the lowest scaling constant that avoids clipping this
particular classification image (see \code{'constant'} scaling below for the formula), so it is not
comparable across classification images with different noise ranges.

\code{'matched'} scaling will match the range of the intensity of the pixels to
the range of the base image pixels. This scaling is nonlinear and depends on the range of both base image
and noise pattern. It is truly suboptimal, because it shifts the 0 point of the noise (that is, pixels that would
not have changed the base image at all before scaling may change the base image after scaling and vice versa). It is
however the quick and dirty way to see how the CI noise affects the base image.

For more control, use \code{'constant'} scaling, where the scaling is independent of
the base image and noise range, but where the choice of constant is arbitrary (provided by the user with
the \code{constant} parameter). The noise is then scale as follows: \code{scaled <- (ci + constant) / (2*constant)}.
Note that pixels can take intensity values between 0 and 1. If your scaled noise exceeds those values,
a warning will be given. You should pick a higher constant (but do so consistently for different classification images
that you want to compare). The higher the constant, the less visible the noise will be in the resulting image.

When creating multiple classification images a good strategy is to find the lowest constant that works for all
classification images. This can be automatized using the \code{autoscale} function.
\end{Details}
%
\begin{Value}
List of pixel matrix of classification noise only, scaled classification noise only, base image only and combined.
\end{Value}
%
\begin{Examples}
\begin{ExampleCode}

# a synthetic square grayscale image stands in for a real base face photo
base_face <- tempfile(fileext = ".png")
png::writePNG(matrix(runif(32 * 32), 32, 32), base_face)

stimulus_path <- tempdir()
generateStimuli2IFC(
  base_face_files = list(face = base_face),
  n_trials = 6,
  img_size = 32,
  stimulus_path = stimulus_path,
  seed = 1,
  ncores = 1,
  nscales = 1,
  save_as_png = FALSE
)
rdata_file <- list.files(stimulus_path, pattern = "\\.Rdata$", full.names = TRUE)[1]

responses <- sample(c(1, -1), 6, replace = TRUE)
ci <- generateCI2IFC(
  stimuli = 1:6, responses = responses, baseimage = "face",
  rdata = rdata_file, save_as_png = FALSE
)

\end{ExampleCode}
\end{Examples}
\inputencoding{utf8}
\HeaderA{generateCINoise}{Generate classification image noise pattern based on set of stimuli (matrix: trials, parameters), responses (vector), and sinusoid}{generateCINoise}
%
\begin{Description}
Generate classification image noise pattern based on set of stimuli (matrix: trials, parameters), responses (vector), and sinusoid
\end{Description}
%
\begin{Usage}
\begin{verbatim}
generateCINoise(stimuli, responses, p)
\end{verbatim}
\end{Usage}
%
\begin{Arguments}
\begin{ldescription}
\item[\code{stimuli}] Matrix with one row per trial, each row containing the 4092 parameters for the original stimulus.

\item[\code{responses}] Vector containing the response to each trial (1 if participant selected original, -1 if participant selected inverted;
this can be changed into a scale).

\item[\code{p}] 3D patch matrix (generated using \code{generateNoisePattern()}).
\end{ldescription}
\end{Arguments}
%
\begin{Value}
The classification image as pixel matrix.
\end{Value}
%
\begin{Examples}
\begin{ExampleCode}
p <- generateNoisePattern(img_size = 32, nscales = 1)
nparams <- max(p$patchIdx)

# two trials, one chosen original (1) and one inverted (-1)
stimuli <- matrix(runif(2 * nparams, -1, 1), nrow = 2)
responses <- c(1, -1)

ci <- generateCINoise(stimuli, responses, p)
\end{ExampleCode}
\end{Examples}
\inputencoding{utf8}
\HeaderA{generateGabor}{Generate single gabor patch}{generateGabor}
%
\begin{Description}
Generate single gabor patch
\end{Description}
%
\begin{Usage}
\begin{verbatim}
generateGabor(img_size, cycles, angle, phase, sigma, contrast)
\end{verbatim}
\end{Usage}
%
\begin{Arguments}
\begin{ldescription}
\item[\code{img\_size}] Integer specifying size of gabor patch in number of pixels.

\item[\code{cycles}] Integer specifying number of cycles the sinusoid should span.

\item[\code{angle}] Value specifying the angle (rotation) of the sinusoid.

\item[\code{phase}] Value specifying phase of sinusoid.

\item[\code{sigma}] of guassian mask on top of sinusoid.

\item[\code{contrast}] Value between -1.0 and 1.0 specifying contrast of sinusoid.
\end{ldescription}
\end{Arguments}
%
\begin{Value}
The sinusoid image with size \code{img\_size}.
\end{Value}
%
\begin{Examples}
\begin{ExampleCode}
generateSinusoid(512, 2, 90, pi/2, 1.0)
\end{ExampleCode}
\end{Examples}
\inputencoding{utf8}
\HeaderA{generateNoiseImage}{Generate single noise image based on parameter vector}{generateNoiseImage}
%
\begin{Description}
Generate single noise image based on parameter vector
\end{Description}
%
\begin{Usage}
\begin{verbatim}
generateNoiseImage(params, p)
\end{verbatim}
\end{Usage}
%
\begin{Arguments}
\begin{ldescription}
\item[\code{params}] Vector with each value specifying the contrast of each patch in noise.

\item[\code{p}] 3D patch matrix (generated using \code{generateNoisePattern()}).
\end{ldescription}
\end{Arguments}
%
\begin{Value}
The noise pattern as pixel matrix.
\end{Value}
%
\begin{Examples}
\begin{ExampleCode}
#params <- rnorm(4092) # generates 4092 normally distributed random values
#s <- generateNoisePattern(img_size=256)
#noise <- generateNoiseImage(params, p)
\end{ExampleCode}
\end{Examples}
\inputencoding{utf8}
\HeaderA{generateNoisePattern}{Generate sinusoid noise pattern}{generateNoisePattern}
%
\begin{Description}
Generate sinusoid noise pattern
\end{Description}
%
\begin{Usage}
\begin{verbatim}
generateNoisePattern(
  img_size = 512,
  nscales = 5,
  noise_type = "sinusoid",
  sigma = 25,
  pre_0.3.0 = FALSE
)
\end{verbatim}
\end{Usage}
%
\begin{Arguments}
\begin{ldescription}
\item[\code{img\_size}] Integer specifying size of the noise pattern in number of pixels.

\item[\code{nscales}] Integer specifying the number of incremental spatial scales. Defaults to 5. Higher numbers will add higher spatial frequency scales.

\item[\code{noise\_type}] String specifying noise pattern type (defaults to \code{sinusoid}; other options: \code{gabor}).

\item[\code{sigma}] Number specifying the sigma of the Gabor patch if noise\_type is set to \code{gabor} (defaults to 25).

\item[\code{pre\_0.3.0}] Boolean specifying whether the noise pattern should be created in a way compatible with older versions of rcicr (< 0.3.0). If you are starting a new project, you should keep this at the default setting (FALSE). There is no reason to set this to TRUE, with the sole exception to recreate behavior of rcicr prior to version 0.3.0.
\end{ldescription}
\end{Arguments}
%
\begin{Value}
List with two elements: the 3D noise matrix with size \code{img\_size}, and an indexing
matrix with the same size to easily change contrasts.
\end{Value}
%
\begin{Examples}
\begin{ExampleCode}
generateNoisePattern(256)
\end{ExampleCode}
\end{Examples}
\inputencoding{utf8}
\HeaderA{generateReferenceDistribution2IFC}{Generates reference distribution}{generateReferenceDistribution2IFC}
%
\begin{Description}
Generates reference distribution of norms for a particular set of task parameters.
\end{Description}
%
\begin{Usage}
\begin{verbatim}
generateReferenceDistribution2IFC(
  rdata,
  iter = 10000,
  ncores = parallel::detectCores() - 1
)
\end{verbatim}
\end{Usage}
%
\begin{Arguments}
\begin{ldescription}
\item[\code{rdata}] String pointing to .RData file that was created when stimuli were generated. This file contains the contrast parameters of all generated stimuli.

\item[\code{iter}] Number of iterations for the simulation (i.e., the number of norms generated with classification images based on random responding).

\item[\code{ncores}] Number of CPU cores to use when re-generating the stimuli (default: detectCores()-1).
\end{ldescription}
\end{Arguments}
%
\begin{Details}
In order to compute the Informational Value metric. Saves its results in the supplied rdata file for later reuse.
\end{Details}
%
\begin{Value}
Nothing. The reference distribution (\code{reference\_norms}) is added to the supplied
\code{rdata} file, so a later call to \code{\LinkA{computeInfoVal2IFC}{computeInfoVal2IFC}} using the same file can
reuse it instead of re-simulating.
\end{Value}
%
\begin{Examples}
\begin{ExampleCode}

# a synthetic square grayscale image stands in for a real base face photo
base_face <- tempfile(fileext = ".png")
png::writePNG(matrix(runif(32 * 32), 32, 32), base_face)

stimulus_path <- tempdir()
generateStimuli2IFC(
  base_face_files = list(face = base_face),
  n_trials = 6,
  img_size = 32,
  stimulus_path = stimulus_path,
  seed = 1,
  ncores = 1,
  nscales = 1,
  save_as_png = FALSE
)
rdata_file <- list.files(stimulus_path, pattern = "\\.Rdata$", full.names = TRUE)[1]

# iter is kept tiny here for a fast example; in practice use iter >= 10000.
# Run from a temp working directory: this function re-generates stimuli via
# generateStimuli2IFC() without forwarding a stimulus_path, so it always
# creates a ./stimuli directory relative to the working directory.
withr::with_dir(tempdir(), {
  suppressWarnings(generateReferenceDistribution2IFC(rdata_file, iter = 3, ncores = 1))
})

\end{ExampleCode}
\end{Examples}
\inputencoding{utf8}
\HeaderA{generateSinusoid}{Generate single sinusoid patch}{generateSinusoid}
%
\begin{Description}
Generate single sinusoid patch
\end{Description}
%
\begin{Usage}
\begin{verbatim}
generateSinusoid(img_size, cycles, angle, phase, contrast)
\end{verbatim}
\end{Usage}
%
\begin{Arguments}
\begin{ldescription}
\item[\code{img\_size}] Integer specifying size of sinusoid patch in number of pixels.

\item[\code{cycles}] Integer specifying number of cycles sinusoid should span.

\item[\code{angle}] Value specifying the angle (rotation) of the sinusoid.

\item[\code{phase}] Value specifying phase of sinusoid.

\item[\code{contrast}] Value between -1.0 and 1.0 specifying contrast of sinusoid.
\end{ldescription}
\end{Arguments}
%
\begin{Value}
The sinusoid image with size \code{img\_size}.
\end{Value}
%
\begin{Examples}
\begin{ExampleCode}
generateSinusoid(512, 2, 90, pi/2, 1.0)
\end{ExampleCode}
\end{Examples}
\inputencoding{utf8}
\HeaderA{generateStimuli2IFC}{Generates 2IFC stimuli}{generateStimuli2IFC}
%
\begin{Description}
Generate stimuli for 2 images forced choice reverse correlation task.
\end{Description}
%
\begin{Usage}
\begin{verbatim}
generateStimuli2IFC(
  base_face_files,
  n_trials = 770,
  img_size = 512,
  stimulus_path = "./stimuli",
  label = "rcic",
  use_same_parameters = TRUE,
  seed = 1,
  maximize_baseimage_contrast = TRUE,
  noise_type = "sinusoid",
  nscales = 5,
  sigma = 25,
  ncores = parallel::detectCores() - 1,
  return_as_dataframe = FALSE,
  save_as_png = TRUE,
  save_rdata = TRUE
)
\end{verbatim}
\end{Usage}
%
\begin{Arguments}
\begin{ldescription}
\item[\code{base\_face\_files}] List containing base face file names used as base images for stimuli. Accepts JPEG and PNG images.

\item[\code{n\_trials}] Number specifying how many trials the task will have (function will generate two images for each trial per base image: original and inverted/negative noise).

\item[\code{img\_size}] Number specifying the number of pixels that the stimulus image will span horizontally and vertically (will be square, so only one integer needed).

\item[\code{stimulus\_path}] Path to save stimuli and .Rdata file to.

\item[\code{label}] Label to prepend to each file for your convenience.

\item[\code{use\_same\_parameters}] Boolean specifying whether for each base image the same set of parameters is used (TRUE) or a unique set is created for each base image (FALSE).

\item[\code{seed}] Integer seeding the random number generator (for reproducibility).

\item[\code{maximize\_baseimage\_contrast}] Boolean specifying whether the pixel values of the base image should be rescaled to maximize its contrast.

\item[\code{noise\_type}] String specifying noise pattern type (defaults to \code{sinusoid}; other options: \code{gabor}).

\item[\code{nscales}] Integer specifying the number of incremental spatial scales. Defaults to 5. Higher numbers will add higher spatial frequency scales.

\item[\code{sigma}] Number specifying the sigma of the Gabor patch if noise\_type is set to \code{gabor} (defaults to 25).

\item[\code{ncores}] Number of CPU cores to use (default: detectCores()-1).

\item[\code{return\_as\_dataframe}] Boolean specifying whether to return a data frame with the raw noise of the stimuli that were generated (default: FALSE). Data frame columns represent pixel values, data frame rows represent stimuli.

\item[\code{save\_as\_png}] Boolean specifying whether to write the stimuli as images to disk (default: TRUE).

\item[\code{save\_rdata}] Boolean specifying whether .RData file with stimulus parameters will be saved (default: TRUE). Note: you always need to save the .RData file so that you can retrieve the stimulus parameters to compute classification images. This function argument exists primarily for internal rcicr use.
\end{ldescription}
\end{Arguments}
%
\begin{Details}
Will save the stimuli as PNGs to a folder, including .Rdata file needed for analysis of data
after data collection. This .Rdata file contains the parameters that were used to generate each stimulus.
\end{Details}
%
\begin{Value}
Nothing, everything is saved to files, unless return\_as\_dataframe is set to TRUE.
\end{Value}
%
\begin{Examples}
\begin{ExampleCode}

# a synthetic square grayscale image stands in for a real base face photo
base_face <- tempfile(fileext = ".png")
png::writePNG(matrix(runif(32 * 32), 32, 32), base_face)

generateStimuli2IFC(
  base_face_files = list(face = base_face),
  n_trials = 4,
  img_size = 32,
  stimulus_path = tempdir(),
  seed = 1,
  ncores = 1,
  nscales = 1
)

\end{ExampleCode}
\end{Examples}
\inputencoding{utf8}
\HeaderA{plotZmap}{Plots a Z-map}{plotZmap}
%
\begin{Description}
Plots a Z-map given a matrix of z-scores that maps onto a specified base image.
\end{Description}
%
\begin{Usage}
\begin{verbatim}
plotZmap(
  zmap,
  bgimage = "",
  sigma,
  threshold = 3,
  mask = NULL,
  decoration = T,
  targetpath = "zmaps",
  filename = "zmap",
  size = 512,
  ...
)
\end{verbatim}
\end{Usage}
%
\begin{Arguments}
\begin{ldescription}
\item[\code{zmap}] A matrix containing z-scores that map onto a given base image. zmap and baseimage must have the same dimensions.

\item[\code{bgimage}] A matrix containing the grayscale image to use as a background. This should be either the base image or the final CI. If not this argument is not given, only the Z-map will be drawn.

\item[\code{sigma}] The sigma of the smoothing that was applied to the CI to create the Z-map.

\item[\code{threshold}] Integer specifying the threshold z-score (default: 3). Z-scores below the threshold will not be plotted on the z-map.

\item[\code{mask}] Optional. A boolean matrix with the same dimensions as zmap. If a cell evaluates to TRUE, the corresponding zmap pixel will be masked. Can also be the filename (as a string) of a black and white PNG image to be converted to a matrix (black = masked).

\item[\code{decoration}] Optional boolean specifying whether the Z-map should be plotted with margins, text (sigma, threshold) and a scale (default: TRUE).

\item[\code{targetpath}] String specifying path to save the Z-map PNG to.

\item[\code{filename}] Optional string to specify a file name for the Z-map PNG.

\item[\code{size}] Integer specifying the width and height of the PNG image (default: 512).

\item[\code{...}] Additional arguments to be passed to raster::plot. Only applied when decoration is TRUE.
\end{ldescription}
\end{Arguments}
%
\begin{Details}
This function takes in a matrix of z-scores (as returned by generateCI) and an Rdata file containing a base image. It returns a Z-map image in PNG format.
Unlisted additional arguments will be passed to raster::plot. For example, a different color palette can be specified using the \code{col} argument. See raster::plot for details.
\end{Details}
%
\begin{Value}
Nothing. It writes a Z-map image.
\end{Value}
%
\begin{Examples}
\begin{ExampleCode}
set.seed(1)
zmap <- matrix(rnorm(64, sd = 5), 8, 8)
plotZmap(zmap, sigma = 3, threshold = 3, decoration = FALSE,
         targetpath = tempdir(), size = 200)
\end{ExampleCode}
\end{Examples}
\inputencoding{utf8}
\HeaderA{simulateNoiseIntensities}{Simulate pixel intensity range for noise}{simulateNoiseIntensities}
%
\begin{Description}
Simulate pixel intensity range for noise
\end{Description}
%
\begin{Usage}
\begin{verbatim}
simulateNoiseIntensities(nrep = 1000, img_size = 512)
\end{verbatim}
\end{Usage}
%
\begin{Arguments}
\begin{ldescription}
\item[\code{nrep}] Number of replications

\item[\code{img\_size}] Size of noise pattern in pixels (one value equal for width and height)
\end{ldescription}
\end{Arguments}
%
\begin{Value}
Matrix with range of noise intensities for each replication
\end{Value}
%
\begin{Note}
This function currently always errors: it references undefined \code{data}/\code{by}
variables when sizing its progress bar (apparently copy-pasted from \code{\LinkA{batchGenerateCI}{batchGenerateCI}}),
and it also ignores the \code{img\_size} argument internally, always simulating at 512px. Not fixed
here; the example below is not run.
\end{Note}
%
\begin{Examples}
\begin{ExampleCode}
## Not run: 
simulateNoiseIntensities(nrep = 10, img_size = 512)

## End(Not run)
\end{ExampleCode}
\end{Examples}
\printindex{}
\end{document}
